\documentclass[11pt]{article}
\usepackage{fontspec}
\setmainfont{Times New Roman}
\setsansfont{Arial}
\setmonofont{Consolas}
\usepackage{amsmath,amssymb,mathtools}
\usepackage{geometry}
\usepackage{microtype}
\geometry{a4paper,margin=2.28cm}
\setlength{\parindent}{1.15em}
\setlength{\parskip}{0pt}
\makeatletter
\renewcommand\footnotesize{\@setfontsize\footnotesize{7.5}{8.5}\selectfont}
\makeatother
\providecommand{\ma}{\mathfrak{a}}
\providecommand{\mb}{\mathfrak{b}}
\providecommand{\mc}{\mathfrak{c}}
\providecommand{\mo}{\mathfrak{o}}
\providecommand{\mr}{\mathfrak{r}}
\emergencystretch=120em
\allowdisplaybreaks
\sloppy
\hbadness=10000
\vbadness=10000
\hfuzz=2pt
\vfuzz=10pt
\raggedbottom
\begin{document}
% Unit: Paper34 section10 v001. Direct Interslavic Latin translation from audited German source lines 564--659, segments P34-S0114--P34-S0128. Source scan pages 18--20 retained as witnesses, and Zenodo record 20822156 was rechecked for this unit. The following Section11 heading is segment P34-S0129 and is intentionally reserved for Section11.
\section*{34. Hiperkompleksne veličiny i teorija predstavjenj}
\emph{Prodolženje: glava II, §§8--14.}

\section*{II. poglavje. Nekomutativna teorija idealov}

\subsection*{§ 10. Razloženje na dvustranne idealy}

Nehaj $\mo$ znovu bude kolco s jediničnym elementom, i nehaj
\begin{equation}
        \mo=\ma_1+\cdots+\ma_n                                  \tag{1}
\end{equation}
bude razloženjem $\mo$ na dvustranne idealy. Togda relacije ortogonalnosti važe takože za idealy:
\begin{equation}
        \ma_i\ma_k=0\quad(i\ne k),
        \qquad
        \ma_i^2=\ma_i.                                             \tag{2}
\end{equation}

\emph{Dokaz.} Nehaj $\mb_1=\ma_2+\cdots+\ma_n$, $\mo=\ma_1+\mb_1$; togda $\ma_1\mb_1\leq \ma_1\cap\mb_1=0$, zato $\ma_1\mb_1=0$, i tym boljše
\[
        \ma_1\ma_i=0\quad(i\ne1);
\]
\[
        \ma_1=\ma_1\mo=\ma_1(\ma_1+\mb_1)=(\ma_1^2,\ma_1\mb_1)=\ma_1^2.
\]

Obratno: \emph{iz (1) i (2) slěduje, že moduly $\ma_i$ sut dvustranne idealy.}
\emph{Dokaz.}
\[
        \mo\ma_i=(\ma_1+\cdots+\ma_n)\ma_i=\ma_i^2=\ma_i,
        \qquad
        \ma_i\mo=\ma_i.
\]

\emph{Prave idealy v kolcu $\ma_i$ sut prave idealy v $\mo$.}
\emph{Dokaz.}
\[
        \mr\mo=\mr(\ma_i+\mb_i)=(\mr\ma_i,\mr\mb_i)=(\mr,0)=\mr.
\]

\emph{Ako $\mr$ jest pravy ideal v $\mo$, togda $\mr$ jest direktna suma pravyh idealov $\mr\ma_i\leq\ma_i$.}
\emph{Dokaz.} $\mr=\mr\mo=(\mr\ma_1,\ldots,\mr\ma_n)$, i $\mr\ma_i\mo=\mr\ma_i$, zato $\mr\ma_i$ sut prave idealy ležeče v $\mr$. Dalje, $\mr\ma_i\leq\ma_i$, zato suma jest direktna:
\[
        \mr=\mr\ma_1+\cdots+\mr\ma_n.
\]

\emph{Osobito, ako $\mr$ jest direktno nerazložimy, togda $\mr$ može iměti samo jednu komponentu; zato $\mr$ mora ležati v někom $\ma_i$.}

Na osnově tyh teoremov vlada se teorijeju idealov v $\mo$, kogda se zna teorija idealov pojedinyh $\ma_i$. Zato se največestěje ograničivamo na dvustranno nerazložime kolca.

\emph{Dva prave idealy, ktore ležet v različnih $\ma_i$, nikogda ne mogut byti operatorno-izomorfne.}
\emph{Dokaz.} Ideal ležeči v $\ma_1$ jest anulovany vsimi drugymi $\ma_k$, ale ne jest anulovany przez $\ma_1$. Naprotiv, ideal ležeči v $\ma_2$ jest anulovany przez $\ma_1$.

Ako tedy možno jest razložiti kolco $\mo$ na dvustranno nerazložime idealy $\ma_1+\cdots+\ma_n$, i vsako $\ma_i$ znovu na jednostranno nerazložime prave idealy $\mr_i'+\mr_i''+\cdots$, togda operatorno-izomorfne $\mr$ treba iskati medžu tymi, ktore prinadležet k tomu samomu $\ma_i$. Ale byvaje, že v tom samom $\ma_i$ jest ješče različne klasy $\mr$ -- t. j. klasy ne operatorno-izomorfne -- kako pokazuje slědujuči primjer.

Nehaj $K$ bude polje racionalnyh čisel,
\[
        \mo=e_1K+e_2K+uK
\]
hiperkompleksny sistem, pri čem $e_i$ i $u$ komutujut s čislami iz $K$, i tablica množenja jest
\[
\begin{array}{c|ccc}
     & e_1&e_2&u\\ \hline
 e_1 & e_1&0&u\\
 e_2 & 0&e_2&0\\
 u   & 0&u&0
\end{array}
\]
Jediničny element jest $e=e_1+e_2$. Elementy $e_i$ izpolnjajut relacije ortogonalnosti. Zato pravo razloženje jest
\[
        \mo=e_1\mo+e_2\mo=(e_1,u)+(e_2),
\]
i lěvo razloženje jest
\[
        \mo=\mo e_1+\mo e_2=(e_1)+(e_2,u).
\]
Poneže $e_2\mo$ i $\mo e_1$ očividno sut nerazložime, $\mo e_2$ i $e_1\mo$ takože moraju byti nerazložime.

Dvustranne razloženje jest nemožno; bo togda jedna komponenta morala by sadržati $e_1\mo$, a druga $e_2\mo$; prva by togda sadržala $u$, ale druga takože $ue_2=u$.

No $e_1\mo$ i $e_2\mo$ ne sut operatorno-izomorfne, poněže imajut različny rang vzhodno k $K$.

\emph{Nehaj $\mo=\ma_1+\cdots+\ma_n$, i nehaj $\ma_i$ sut dvustranno nerazložime. Togda one sut jednoznačno opreděljene.}

\emph{Dokaz.} Nehaj kromě togo $\mo=\mc_1+\cdots+\mc_n$, tedy
\[
        \ma_i=\ma_i\mo=(\ma_i\mc_1,\ldots,\ma_i\mc_n).
\]
Poslědnja suma jest direktna, poněže $\ma_i\mc_k\leq\mc_k$ i $\leq\ma_i$; ale poněže $\ma_i$ sut nerazložime, vsi $\ma_i\mc_k$ moraju byti nuljevy kromě jednogo: $\ma_i\mc_{j_i}$. Zato
\[
        \ma_i=\ma_i\mc_{j_i}.
\]
Takože:
\[
        \mc_{j_i}=\mo\mc_{j_i}=\ma_1\mc_{j_i}+\cdots+\ma_n\mc_{j_i},
\]
i $\ma_i\mc_{j_i}\ne0$, zato vsi ostale sut $=0$,
\[
        \mc_{j_i}=\ma_i\mc_{j_i}=\ma_i,
\]
čto bylo dokazati.
\end{document}
